\documentclass[10pt,a4paper]{amsart}
\usepackage[margin=19mm]{geometry}
\usepackage{amsmath,amssymb,mathtools}
\usepackage[scr=rsfs,cal=euler]{mathalfa}
\usepackage{xfrac}
\usepackage[unicode,hidelinks,bookmarks=false]{hyperref}
\newcommand{\ie}{i.e.\ }
\newcommand{\cf}{cf.\ }
\newcommand{\resp}{respectively\ }
\newcommand{\Subsection}{\subsection}
\providecommand{\nobreakdash}{\nobreak\hbox{-}}
\setlength{\emergencystretch}{2em}
\multlinegap0pt
\usepackage{xcolor}
\usepackage{bm}
\usepackage{amssymb}
\usepackage{tikz}
\usepackage{enumerate}
\usepackage[shortlabels]{enumitem}
\usepackage{mathtools}
\newtheorem{lemma}{Lemma}
\newcommand{\E}{\mathbf{E}}
\newcommand{\eps}{\varepsilon}
\newcommand{\prob}{\mathbf{P}}
\title{Inhomogeneous random graphs \\ with infinite-mean fitness variables}
\author{Luca Avena, Diego Garlaschelli, Rajat Subhra Hazra, Margherita Lalli}
\date{Source: arXiv:2212.08462v3 (CC BY 4.0)}
\begin{document}
\maketitle

\noindent\emph{Source and licence.} Inhomogeneous random graphs \\ with infinite-mean fitness variables. Version arXiv:2212.08462v3 (CC BY 4.0), distributed by arXiv under the Creative Commons Attribution 4.0 International licence, http://creativecommons.org/licenses/by/4.0/, as recorded in the arXiv licence metadata for this exact version and shown on the abstract page https://arxiv.org/abs/2212.08462v3. Full licence notice: \texttt{licenses/arxiv-2212.08462-CC-BY-4.0.txt}.\par\smallskip

\noindent\textit{Editorial scope.} Counted unit: the article's lemma cond:degree (source0.tex lines 1199--1207). The article's complete original proof of it is delivered with it (source0.tex lines 1208--1236, one \texttt{proof} environment of 29 lines). No mathematics is written, repaired or re-derived: the statement and the proof are the article's own lines, in the article's own notation, and no retained line is substituted or re-flowed.  Retained uncounted as the source-local context the proof needs: lines 1172--1176; lines 1190--1193.  Nothing was omitted from any retained span, so the package carries no omission record.  No retained line contains a citation, so the wrapper carries no bibliography and no bibliography entry.  No retained line contains a figure environment, an image-inclusion command or drawing code, so no image is delivered and the package declares no asset dependency.  Editorial formatting only, in addition: an AMS \texttt{amsart} preamble that restates the article's own declarations and macros used by the retained lines, namely the environments \texttt{lemma} on separate counters, together with the standard packages those lines need.  The retained environments print in this document as Lemma 1 in order of appearance, whereas the article prints the same items with the numbering of its own full document.  The complete native source of the article as distributed in the arXiv e-print for this exact version is bundled unchanged as \texttt{sources/134976/source0.tex}.
\begin{lemma}\label{lemma:incompletegamma}
Let $s=(1-\alpha)\in (0,1)$ and let
$$\Gamma(s, x)= \int_x^\infty e^{-z}z^{s-1}dz, \, \, x>0.$$ Then
$$e^{-x} (1+x)^{s-1}\le \Gamma(s,x) \le e^{-x}\Gamma(s) \left(1+\frac{x}{s}\right)^{s-1}.$$
\end{lemma}

\begin{lemma}\label{lemma:surprising}
For $\alpha\in [0,1]$ and $x>0$, we have
$$1-e^{-x} \le x^{\alpha}.$$
\end{lemma}

\begin{lemma} 
 For $i\in [n]$ and $j\neq i$ and $w>1$, we have
\begin{equation}\label{cond:degree}
 \eps_n^{\alpha} w^{\alpha} e^{-(1+\alpha) \eps_n w}\le \E[ p_{ij}\mid W_i=w] \le \eps_n^{\alpha}w^\alpha(1+\Gamma(1-\alpha) e^{-\eps_n w})
 \end{equation}
 In particular, for fixed $w>1$, 
 
 \begin{equation}\label{asymp}\E[ p_{ij}\mid W_i=w]\sim \Gamma(1-\alpha)\eps_n^{\alpha} w^{\alpha}, \text{ as $\eps_n\downarrow 0$}.\end{equation}
\end{lemma}

\begin{proof}
We drop the subscript $n$ from $\eps_n$ for notational simplicity.
\begin{align*}
    \E[ p_{ij}\mid W_i=w]&=  \E\left[ \left(1-e^{-\eps W_j w}\right)\right]\\
    &= \int_0^\infty (1-e^{-\eps w u}) F_W(du)= \int_0^\infty \int_0^u e^{-\eps w z} (\eps w) dz F_W(du)\\
    &= \eps w \int_0^\infty e^{-\eps w z}\prob(W_j>z) dz.
\end{align*}
where we used Fubini's theorem to swap the integrals. Using the fact that $W_j$ are supported on $(1,\infty)$ we get that
\begin{align}
    \E[ p_{ij}\mid W_i=w]&=\eps w \int_0^1 e^{-\eps w z} dz+  \eps w\int_1^\infty e^{-\eps w z}z^{-\alpha} dz\nonumber\\
    &=(1-e^{-\eps w})+  \eps^\alpha w^\alpha \int_{\eps w}^\infty e^{-z}z^{-\alpha} dz=  (1-e^{-\eps w}) + \eps^\alpha w^\alpha \Gamma(1-\alpha, \eps w), \label{eq:intermediate}
\end{align}
where $\Gamma(1-\alpha, \eps w)$ is an incomplete Gamma function. Now we can use the bounds from Lemma \ref{lemma:incompletegamma} and Lemma \ref{lemma:surprising}. So for the upper bound we have
$$ \E[ p_{ij}\mid W_i=w]\le \eps^{\alpha}w^{\alpha}+ \eps^\alpha w^\alpha \Gamma(1-\alpha) e^{-\eps w}\left( 1+ \frac{\eps w}{1-\alpha}\right)^{-\alpha}$$
Now using the trivial bounds  $1+ \frac{\eps w}{1-\alpha}\ge 1$ we have
$$\E[ p_{ij}\mid W_i=w]\le \eps^{\alpha} w^{\alpha}(1+\Gamma(1-\alpha) e^{-\eps w}).$$

For the lower bound, we ignore the first term in \eqref{eq:intermediate} and use the lower bound in Lemma \ref{lemma:incompletegamma}. We have 
\begin{align*}
  \E[ p_{ij}\mid W_i=w]\ge \eps^\alpha w^{\alpha}\Gamma(1-\alpha, \eps w)
  \ge \eps^{\alpha} w^\alpha e^{-\eps w}(1+\eps w)^{-\alpha}\ge \eps^{\alpha}  w^{\alpha} e^{-\eps w}e^{-\alpha \eps w}
\end{align*}
%\begin{align*}
 % \E[ p_{ij}\mid W_i=w]&\ge \eps^\alpha w^{\alpha}\Gamma(1-\alpha, \eps w) \\
  %&\ge \eps^{\alpha} w^\alpha e^{-\eps w}(1+\eps w)^{-\alpha}\\
  %&\ge \eps^{\alpha}  w^{\alpha} e^{-\eps w}e^{-\alpha \eps w}
%\end{align*}
where we used $1+x\le e^x$, for $x>0$ and hence $(1+x)^{-\alpha}\ge e^{-\alpha x}$. This completes the proof of the lower bound. The asymptotic estimate in \eqref{asymp} follows from \eqref{eq:intermediate}.
\end{proof}

\end{document}
